\chapter{Actuaciones reales, gravedad y resistencia}\label{ch:27}
\objetivos{Formular una trayectoria puntual con marco y ángulos definidos; interpretar pérdidas gravitatorias, aerodinámicas y de dirección; comprobar una integración mediante balances y casos límite.}
\section{Ecuación vectorial del vehículo}
En un marco inercial y con un modelo de masa puntual,
\begin{equation}\label{eq:trajectoryvector}
 m\dot{\mathbf v}=\mathbf F_T+\mathbf F_D+\mathbf F_L
 -m\frac{\mu}{r^3}\mathbf r,\qquad
 \dot{\mathbf r}=\mathbf v,\qquad \dot m=-\dot m_p.
\end{equation}
La fuerza aerodinámica depende de velocidad relativa al aire, que puede diferir de la inercial por rotación atmosférica y viento. En un marco solidario con un planeta aparecerían términos de Coriolis y centrífugo. No se mezclan ecuaciones de un marco con velocidades iniciales de otro.

La gravedad local es $g(r)=\mu/r^2$, distinta del valor convencional $g_0$ utilizado para definir impulso específico. Un cambio de altura modifica gravedad física, pero no la conversión $c=g_0I_{sp}$.
\section{Modelo plano sobre un cuerpo esférico}
Para movimiento en un plano sin rotación, sean $v$ rapidez inercial, $r$ radio, $\gamma_f$ ángulo de trayectoria sobre la horizontal local y $\alpha_T$ inclinación del empuje respecto a velocidad. Con resistencia opuesta al movimiento y sustentación positiva hacia aumento de $\gamma_f$,
\begin{align}
 \dot r&=v\sin\gamma_f,\\
 \dot\theta&=\frac{v\cos\gamma_f}{r},\\
 \dot v&=\frac{F\cos\alpha_T-D}{m}-\frac{\mu}{r^2}\sin\gamma_f,\label{eq:vdot}\\
 \dot\gamma_f&=\frac{F\sin\alpha_T+L}{mv}
 +\left(\frac{v}{r}-\frac{\mu}{vr^2}\right)\cos\gamma_f.
\end{align}
El término $v/r$ representa el giro de la horizontal local. El sistema es singular en $v=0$ por la parametrización angular, no porque la dinámica vectorial deje de existir. Cerca del inicio conviene usar componentes cartesianas o condiciones que eviten dividir por cero.\cite{MIT32,MIT33}

Estas ecuaciones no describen por sí solas actitud del cuerpo ni flexibilidad. El ángulo de empuje puede ser una entrada de un modelo puntual, pero realizarlo exige dinámica y autoridad de control adicionales.
\section{Pérdidas como descomposición de un balance}
Integrando \eqref{eq:vdot},
\begin{align}
 v_f-v_0={}&\int\frac{F}{m}dt
 -\int\frac{F(1-\cos\alpha_T)}{m}dt\\
 &-\int\frac{D}{m}dt-\int g(r)\sin\gamma_f\,dt.
\end{align}
La primera integral es delta-v propulsivo escalar. La segunda cuantifica la parte que no aumenta rapidez por inclinación del empuje respecto al movimiento en esta descomposición. Las dos últimas son contribuciones de resistencia y gravedad proyectada. Se debe indicar el intervalo y convenio antes de darles el nombre de pérdidas.

Durante un descenso o maniobra orbital, la contribución gravitatoria proyectada puede cambiar de signo. La gravedad no destruye energía mecánica en un campo conservativo; intercambia cinética y potencial. El término pérdida gravitatoria expresa una comparación propulsiva de trayectoria, no disipación termodinámica.
\section{Ascenso vertical elemental}
Para empuje vertical, gravedad constante, sin resistencia y velocidad efectiva constante,
\[
 \dot v=-c\frac{\dot m}{m}-g.
\]
La integración hasta tiempo $t_b$ da
\begin{equation}\label{eq:vertical}
 v_b-v_0=c\ln\frac{m_0}{m_b}-gt_b.
\end{equation}
El término $gt_b$ puede ser importante aunque la relación de masas sea favorable. A igual capacidad ideal, una duración más larga aumenta esta contribución en el modelo vertical. Esa observación no implica que reducir duración sin límites sea un diseño realizable: aceleración, estructura y potencia imponen restricciones.

Si $F<mg$ al inicio y no hay soporte ni velocidad ascendente suficiente, no se deduce un despegue vertical. El cociente empuje-peso inicial es una comprobación previa. Una integración que permita al vehículo atravesar el suelo sin un modelo de contacto puede producir números sin representar el escenario planteado.
\section{Resistencia, presión dinámica y atmósfera}
La resistencia se representa por
\[
 D=\frac12\rho_{atm}V_{rel}^2C_D S.
\]
El coeficiente depende, entre otras variables, de Mach, Reynolds y configuración. La presión dinámica es $q=\rho_{atm}V_{rel}^2/2$, distinta de presión estática. Su máximo puede aparecer a una altura intermedia porque la velocidad crece mientras la densidad disminuye.

Un modelo exponencial $\rho=\rho_0e^{-h/H}$ es una aproximación de escala con altura característica constante. No es una atmósfera estándar completa ni reproduce variación térmica y composición a cualquier altura. En problemas originales se declarará cuándo se utiliza para estudiar tendencias, evitando atribuirle precisión meteorológica.
\section{Energía y momento angular como controles}
En un tramo balístico sin resistencia ni empuje, la energía específica
\[
 \mathcal E=\frac{v^2}{2}-\frac{\mu}{r}
\]
y el momento angular específico $\mathbf h=\mathbf r\times\mathbf v$ permanecen constantes. Una integración numérica puede comprobar ambos para un caso de referencia. Disminuir tolerancia o paso debe reducir errores de discretización hasta el límite de redondeo y modelo.

Con empuje y resistencia,
\[
 \frac{d\mathcal E}{dt}=\frac{\mathbf F_T\cdot\mathbf v+\mathbf F_D\cdot\mathbf v+\mathbf F_L\cdot\mathbf v}{m}
\]
en el campo central fijo. Para sustentación perpendicular a velocidad inercial su trabajo es cero; con atmósfera móvil, la perpendicularidad se refiere a velocidad relativa y el balance debe tratarse con cuidado. Los controles se adaptan al marco y al modelo, no se copian indiscriminadamente.
\section{Qué demuestra una simulación docente}
Una trayectoria calculada con masas, atmósfera y coeficientes idealizados demuestra el comportamiento de esas ecuaciones. Para predecir un lanzamiento real se necesitan datos de vehículo, aerodinámica, dispersión, control y restricciones. Una integración convergida puede resolver con precisión un modelo insuficiente.

La comprobación del manual comparará el caso vertical con solución analítica, cerrará masa y revisará energía en tramos sin fuerzas no conservativas. La validación con telemetría real queda fuera del alcance y se mantiene pendiente.
\figura{c27}{Solución de ecuación vertical con gravedad constante, sin suelo ni resistencia. El empuje inicial supera el peso; el modelo no incluye atmósfera ni variación de gravedad.}
\section{Ejercicios propuestos}
\begin{enumerate}
\item Derive \eqref{eq:vertical} y determine el efecto de duplicar duración a igual relación de masas.
\item Explique por qué $gt$ no es la pérdida gravitatoria general de cualquier trayectoria.
\item Distinga presión estática y dinámica, y razone la posibilidad de un máximo de $q$ durante ascenso.
\item Proponga dos invariantes para verificar un integrador de un tramo orbital sin empuje ni resistencia.
\end{enumerate}
\resumen{Las actuaciones reales dependen de historia temporal, dirección y fuerzas externas. Las pérdidas son términos de un balance definido, y gravedad no equivale a disipación. La convergencia numérica y los invariantes comprueban una integración, pero no validan automáticamente el modelo de lanzamiento.}
